\section{Minimal teaching information}\label{sec:minimal-information}
The student need not recover the entire normal vector in
Eq.~\eqref{eq:cost-preserving-loss}. We characterize the part affecting
its cost comparisons under a specified scalar-feedback interface.

\subsection{Interface}
At a fixed teacher state, let
\begin{equation}\label{eq:minimal-model}
 Q(u)=a\|u\|^2+b^\top u+c,\qquad
 v=\arg\min_{u\in U}Q(u),\qquad n=2av+b\in-N_U(v).
\end{equation}
Here $a>0$, $U$ is a centered ball or box of radius or half-width $R>0$, and
$N_U(v)=\{w:w^\top(u-v)\leq0\text{ for all }u\in U\}$.
The receiver knows $a,U,v$ and orthonormal $P\in\R^{d\times m}$;
its actions lie in $S\cap U$, $S=\operatorname{range}(P)$.
It has no access to $b,n$. A feasible query returns $Q(u)-Q(0)$, giving
\begin{equation}\label{eq:minimal-response}
 y(u)=Q(u)-Q(0)-a\|u\|^2+2av^\top u=n^\top u.
\end{equation}
We count additional scalar responses, with target and baseline supplied.
The bound excludes arbitrary encoded messages, derivative access and
learners allowed to recompute $b$ from the full analytic model.
This is an abstract restricted teaching interface, not a demonstrated
communication bottleneck in a deployed controller.

Let orthonormal $\Gamma\in\R^{d\times k}$ span
$K=\operatorname{span}(-N_U(v))$; $k=0$ for interior targets.
For student actions, only $\beta=P^\top n$ matters:
\begin{equation}\label{eq:student-normal}
 Q(Pz)-Q(Pz')=a(\|z\|^2-\|z'\|^2)
       +(-2aP^\top v+\beta)^\top(z-z').
\end{equation}
Normal differences in $K\cap S^\perp$ are therefore irrelevant.
The recovery requirement covers every feasible pair in $S\cap U$;
it is stronger than recovering only the comparisons used by a
particular fitted network.

\begin{proposition}\label{prop:minimal-information}
Fixed feasible queries with row matrix $W$ determine every cost
difference in $S\cap U$, uniformly over models sharing $a,U,v$, iff
\begin{equation}\label{eq:kernel-information}
 \ker(W\Gamma)\subseteq\ker(P^\top\Gamma).
\end{equation}
The minimum worst-case count is
\begin{equation}\label{eq:information-rank}
 r=\operatorname{rank}(P^\top\Gamma)\leq\min(k,m).
\end{equation}
This lower bound includes deterministic adaptive queries and is attained
by a feasible nonadaptive construction.
\end{proposition}
\begin{proof}
Write $n=\Gamma\eta$. Observations $W\Gamma\eta$ determine
$P^\top\Gamma\eta$ precisely under Eq.~\eqref{eq:kernel-information}.
Since $S\cap U$ contains a neighborhood of zero in $S$, a coefficient
difference changes a feasible cost difference. Necessity holds within
the normal cone: small opposite perturbations of a relative-interior
normal along an unobserved direction in $K$ remain admissible and,
by strong convexity, preserve the unique target $v$.

With $q<r$ queries, the row space of $W\Gamma$ cannot contain that
of $P^\top\Gamma$. For adaptive queries, follow the transcript at a
relative-interior normal and perturb within the common kernel of those
queries. Both perturbations preserve every response and subsequent
query, but change the student coefficient.

For attainment, choose orthonormal $E\in\R^{m\times r}$ spanning
$\operatorname{range}(P^\top\Gamma)$. Query $u_j=\alpha PE_j$,
$\alpha=R/2$, where $R$ is the ball radius or box half-width. Each
query has norm $\alpha$ and is feasible. The responses satisfy
$y=\alpha E^\top\beta$, so $\beta=Ey/\alpha$. If $r=0$, no query
is required.
\end{proof}

\subsection{Direct-query controls}
For any orthogonal $H\in\R^{m\times m}$, the $m$ feasible queries
$u_j=\alpha PH_j$ give $y=\alpha H^\top\beta$ and
$\beta=Hy/\alpha$. Thus both the student coordinate basis and a random
orthogonal student basis recover exactly. A hybrid choosing full-normal
recovery when $k<m$ and student-basis recovery otherwise uses
$\min(k,m)$ responses. It already attains the bound whenever
$r=\min(k,m)$; relevant queries then offer no count advantage.
Interior targets need no query for any method.

A strict query-count advantage over the hybrid occurs when $r<\min(k,m)$.
We focus on the nontrivial case $r>0$. For example,
let $d=8$, $K=\operatorname{span}(e_1,e_2,e_3)$ and
\begin{equation}\label{eq:shared-actuation}
 P=[g,e_4,e_5,e_6]H,\qquad
 g=(e_1+e_2+e_3)/\sqrt3,
\end{equation}
with a generic orthogonal $4\times4$ matrix $H$. Here $k=3,m=4,r=1$.
All three active coordinate normals and all four student coordinate
directions are visible, so deleting zero directions alone does not
attain the bound. The student can act on only their shared combination.
Recognizing this dependency, by any equivalent rank-revealing method,
reduces three hybrid queries to one.

\subsection{Learning rule}
At each cached state, form $\Gamma,E$ from $U,v,P$, collect the
$r$ responses, and train $u=Pz_\theta(x)$ with
\begin{equation}\label{eq:minimal-loss}
 L_{\rm relevant}(z;x)=a\|z\|^2+
          (-2aP^\top v+Ey/\alpha)^\top z=Q(Pz)-Q(0).
\end{equation}
This preserves the full-information objective and its parameter
gradients where differentiable, without reconstructing $n$ or
requiring the student to realize $v$. Response error $e$ gives
\begin{equation}\label{eq:minimal-noise}
 \|\widehat\beta-\beta\|\leq\frac{\|e\|}{\alpha},\qquad
 |\widehat{\Delta Q}-\Delta Q|
 \leq\frac{\|e\|}{\alpha}\|z-z'\|.
\end{equation}
The result is statewise and requires known curvature, exact geometry
and fixed $P$; ordinary hidden width does not imply a fixed output
space. Equation~\eqref{eq:minimal-noise} is a local error bound.

\subsection{Approximate alignment}
Exact rank can increase under arbitrarily small perturbations of $P$.
If a bound $\|n\|\leq M$ is available, a truncated construction instead
gives a continuous error tradeoff. Let $E_q$ contain the first $q$ left
singular vectors of $P^\top\Gamma$, padding its singular values $\sigma_j$
with zeros beyond the spectrum.
Querying $\alpha PE_q$ and decoding noisy responses as above yields
\begin{equation}\label{eq:truncated-feedback}
 \|\widehat\beta-\beta\|^2\leq
 M^2\sigma_{q+1}^2+\|e\|^2/\alpha^2.
\end{equation}
Indeed, $\beta=P^\top\Gamma\eta$, $\|\eta\|=\|n\|\leq M$;
the omitted singular components have norm at most $M\sigma_{q+1}$
and are orthogonal to the decoded response error. Multiplying the
square root of this bound by $\|z-z'\|$ bounds the cost-difference
error. This is an approximation upper bound requiring $M$, not an
exact recovery or approximate-query optimality claim. Numerical rank
tolerances must therefore be interpreted against the discarded spectrum.

\subsection{Ball and box constraints}
At an active ball target, $n=-\lambda v$, $k=1$. If $P^\top v\ne0$,
one scalar is necessary and sufficient. The target gain gives
\begin{equation}\label{eq:gain-normal}
 G=Q(0)-Q(v)=(a+\lambda)R^2,\qquad\lambda=G/R^2-a.
\end{equation}
It recovers the full radial-slack loss, whereas scalar weighting of
MSE omits that slack. If $P^\top v=0$, no additional information is needed.

For a box with $k$ active coordinates, $\Gamma$ contains their coordinate
vectors and generically $r=\min(k,m)$. Gain alone can fail: with
$a=1$, $U=[-1,1]^2$, $b_A=(-3,-5)$ and $b_B=(-5,-3)$ both yield
$v=(1,1)$ and $G=6$. Yet costs relative to zero at $(1,0),(0,1)$
are $(-2,-4)$ for $A$ and $(-4,-2)$ for $B$.

In the HJ model, $b$ contains $B^\top\nabla V^\phi$; the rank thus
counts hidden future-cost directions relevant to this student.
For an approximate continuation model, recovery preserves that model;
policy transfer still requires the stated remainder and occupancy bounds.
